\section{Introduction}

Gigalis is the public-interest group responsible for digital services in the Pays de la Loire region. Among its roles it acts as a central purchasing body: it builds pooled framework agreements in cloud, cybersecurity, networks and artificial intelligence, which its members may use but are never obliged to. The strategic value of those agreements therefore depends on Gigalis anticipating its members' purchasing needs rather than observing them after the fact. The question put to this internship was: using historical data on digital public procurement, can we estimate the probability that a contract or a technology segment generates an identifiable purchasing need within the next twelve to twenty-four months?

The work was carried out entirely on open data --- the notices published in the \emph{Bulletin officiel des annonces des marchés publics} (BOAMP), 1,620,712 of them between 2015 and 2025, reduced to a study cohort of 3,800 awarded digital procurements in the Grand Ouest. The means were those of a standard workstation: Python, scikit-learn, lifelines, statsmodels and ruptures, with an end-to-end reproducible pipeline.

It became clear early on that the main difficulty was not estimating a duration but defining what is being measured. The BOAMP carries no field indicating that one contract renews another. The outcome variable therefore had to be constructed, and that construction governs everything downstream. This report presents the construction, the survival analysis it enables, the technology segmentation learnt from the text of the notices, the descriptive trend analysis, and a critical reading of what these results support --- and of what they do not.
